\makereporttitle{2024}{mai}

\makereportpreamble

\textbf{Activitate de cercetare:}

În această lună am continuat să studiez articolul „\emph{Calculating eigenvalues and eigenvectors of parameter-dependent hamiltonians using an adaptative wave operator method}” a lui Arnaud Leclerc și Georges Jolicard, publicat în 2020 în \emph{Journal of Chemical Physics}, volumul 152, numărul 20.

Am luat legătura cu autorii lucrării și am obținut codul sursă comentat, scris în MATLAB, care implementează atât algoritmul \emph{adaptive wave operator}, cât și algoritmii \emph{standard wave operator} sau \emph{Davidson} pentru comparație. Am început să mă familiarizez cu codul și să studiez cum ar putea fi adaptat acesta pentru a studia Hamiltonienii care depind de un parametru din mecanica cuantică.

\vspace{0.5em}

\textbf{Cursuri și seminare științifice:}

Luna aceasta am participat la două cursuri doctorale:
\begin{itemize}
    \item Cursul de ,,Quantum Information Theory'', ținut de prof.~dr.~Aurelian Isar, directorul Departamentului de Fizică Teoretică al IFIN-HH;

    \item Cursul de ,,Fizica particulelor elementare: SM și MSSM'', ținut de C.S.~I.~Călin Alexa, membru al DFT și CERN;
\end{itemize}

\vspace{0.5em}

\textbf{Observațiile conducătorului științific}:

\vspace{3em}

\makereportpostamble{30 mai 2024}
